% ==============================================================================
% Section 2: Perceptual Foundations
% ==============================================================================
% VERSION: 2.0 (Comprehensive Redraft - 30 Dec 2025)
% Research Integration: Hong (2024), Nölle (2012), Braun (2017), Sekulovski (2007)
% Target Length: 3,000–3,500 words
% Dependencies: None (foundational section)
% ==============================================================================

\section{Perceptual Foundations}
\label{sec:perceptual-foundations}

The design of temporally coherent color progressions requires a rigorous foundation in perceptual color science. This section establishes three critical principles that underpin all subsequent design decisions: (1) color space is a Riemannian manifold, not a Euclidean vector space; (2) perceptual uniformity in three dimensions is mathematically impossible; and (3) temporal color perception differs fundamentally from spatial color perception. These are not engineering conveniences—they are empirical facts with cascading implications for how color journeys must be constructed.

% ------------------------------------------------------------------------------
% 2.1 The Riemannian Structure of Color Space
% ------------------------------------------------------------------------------
\subsection{Color Space as Riemannian Manifold}
\label{sec:riemannian-color}

\subsubsection{Theoretical Framework}

Recent comprehensive measurements of human color discrimination thresholds confirm that color space must be modeled as a \emph{Riemannian manifold}—a space that is locally Euclidean but globally curved \citep{hong2024}. This finding has profound implications for color interpolation and distance measurement.

In a Euclidean space, the shortest path between two points is a straight line, and distances can be computed using Pythagorean geometry. In a Riemannian manifold, the shortest path is a \emph{geodesic}—a curve that follows the intrinsic geometry of the space \citep{hong2024}. The perceptual distance between two colors corresponds to the geodesic distance: the accumulated perceptual thresholds along the path of minimal discrimination \citep{hong2024}.

Hong et al.\ \citeyearpar{hong2024} characterized this structure empirically through approximately 6,000 discrimination threshold trials per participant, mapping elliptical contours of equal perceptual distance across color space. Their findings reveal systematic patterns:

\begin{enumerate}
    \item \textbf{Radial orientation:} Major axes of threshold ellipses are consistently oriented toward the achromatic center \citep{hong2024}
    \item \textbf{Local metric variation:} The perceptual ``size'' of a just-noticeable difference varies by location in color space
    \item \textbf{Geodesic curvature:} Perceptually optimal paths between distant colors are curved, not straight, in standard colorimetric coordinates
\end{enumerate}

This Riemannian framework has historical foundations extending to Schrödinger's extension of Helmholtz's line element \citep{roberti2023}. Roberti and Peruzzi \citeyearpar{roberti2023} note that Schrödinger distinguished between \emph{lower color metrics} (where affine geometry suffices for color mixture) and \emph{higher color metrics} (where Riemannian geometry is necessary to translate color differences into perceptual distances). The engine operates in the domain of higher color metrics.

\subsubsection{Practical Implications}

For color journey construction, the Riemannian structure has three critical consequences:

\begin{description}
    \item[Linear interpolation in RGB is perceptually nonuniform] Interpolating between saturated blue $(0, 0, 255)$ and saturated yellow $(255, 255, 0)$ in sRGB produces a midpoint near $(127, 127, 127)$—a desaturated grey—because the straight line passes through the achromatic region. The perceptually smooth path would curve around the exterior of color space.
    
    \item[Distance metrics must respect local geometry] A color difference $\Delta E = 5.0$ has different perceptual magnitudes in different regions. Near the achromatic center, where discrimination thresholds are smaller, this represents a larger perceptual step than at high chroma.
    
    \item[Gamut mapping must follow geodesics] When projecting out-of-gamut colors to displayable coordinates, the projection should follow perceptual geodesics rather than straight lines in RGB (\S\ref{sec:gamut-management}). Hong et al.'s finding that discrimination ellipses orient radially toward the achromatic center \citep{hong2024} justifies reducing saturation (moving toward achromatic) as the perceptually minimal modification.
\end{description}

\subsubsection{Computational Approximation}

True geodesic computation in a Riemannian manifold requires solving differential equations or variational problems—computationally expensive for real-time applications. The engine adopts a pragmatic approach: work in a color space where \emph{Euclidean straight lines approximate geodesics sufficiently well}. This is the definition of perceptual uniformity.

Modern perceptually uniform color spaces (OKLab, CAM16-UCS) are designed such that Euclidean distance $\Delta E$ in the coordinate representation approximates geodesic distance in the underlying perceptual manifold \citep{ottosson2020,safdar2017}. The approximation is not perfect—local curvature remains \citep{hong2024}—but the error is small enough for practical interpolation tasks.

The engine thus achieves a balance: acknowledging the Riemannian nature of color space theoretically while exploiting Euclidean approximations computationally. All subsequent design parameters ($\Delta_{\min}$, $\Delta_{\max}$, velocity weights) are interpreted as \emph{locally Euclidean} values in OKLab coordinates.

% ------------------------------------------------------------------------------
% 2.2 The Impossibility of 3D Euclidean Uniformity
% ------------------------------------------------------------------------------
\subsection{The 720° Topology and Euclidean Impossibility}
\label{sec:topology-impossibility}

\subsubsection{Mathematical Statement}

A fundamental result in perceptual color theory is that no three-dimensional Euclidean space can achieve perfect perceptual uniformity \citep{nolle2012}. This is not a limitation of current measurement techniques—it is a topological impossibility arising from the ``super-importance of hue'' first identified by Judd \citep{judd1940}.

Nölle et al.\ \citeyearpar{nolle2012} provide a rigorous derivation. The perceptual distance around a constant-lightness, constant-saturation hue circle is given by the line integral:

\begin{equation}
    U(S) = \int_{0}^{2\pi} \sqrt{g_{HH}(S)} \, dH
    \label{eq:hue-circumference}
\end{equation}

\noindent where $g_{HH}$ is the metric tensor component for hue angle $H$ at saturation $S$. Empirical measurements of color discrimination thresholds yield \citep{nolle2012}:

\begin{equation}
    U(S=1) \approx 12.65 \approx 4\pi
    \label{eq:720-result}
\end{equation}

In a Euclidean space, the circumference of a unit circle is $2\pi$. Color space requires approximately \emph{twice} this value—the hue circle has an effective ``circumference'' of 720° rather than 360° \citep{nolle2012}.

\subsubsection{Physical Interpretation}

What does this mean physically? Consider walking around a hue circle at constant lightness and saturation, taking steps of one just-noticeable difference (JND). In a Euclidean space, you would return to your starting point after $2\pi r$ JND steps, where $r$ is the radius. In color space, you require $4\pi r$ steps—\emph{double} the expected number.

This ``extra length'' cannot be accommodated in three-dimensional Euclidean geometry. It implies that color space has intrinsic curvature: the space is \emph{non-Euclidean} \citep{nolle2012,roberti2023}. Any attempt to embed color perception in a flat 3D coordinate system must introduce distortions—either uneven discrimination thresholds (as in CIELAB) or hue-dependent perceptual compression (as in HSV/HSL).

\subsubsection{Implications for Existing Color Spaces}

This result explains why even carefully designed spaces like CIELAB exhibit perceptual nonuniformities \citep{kong2021}:

\begin{itemize}
    \item \textbf{CIELAB (1976):} Achieves approximate uniformity in spatial discrimination tasks but fails for temporal transitions. Kong \citeyearpar{kong2021} notes explicitly: ``CIELAB is not a useful space to predict the perception of dynamic colored light. Today, no color spaces are available that accurately predict the visibility of color differences over time.''
    
    \item \textbf{CAM16-UCS:} Provides excellent uniformity but at high computational cost (iterative convergence) \citep{safdar2017}
    
    \item \textbf{HSV/HSL:} Euclidean spaces with poor perceptual correspondence; hue angle does not represent equal perceptual steps
\end{itemize}

The engine adopts OKLab as a pragmatic compromise: a space that \emph{minimizes} perceptual distortion while remaining computationally tractable (\S\ref{sec:oklab-implementation}). The 720° result remains relevant for understanding why perfect uniformity is unattainable and why certain design choices (e.g., velocity weights, closure thresholds) require empirical calibration rather than pure geometric derivation.

\subsubsection{Connection to Möbius Loop Topology}

The 720° result has a surprising implication for loop closure strategies (\S\ref{sec:loop-strategies}). A Möbius strip—a surface with a single half-twist—also requires traversing 720° to return to the original orientation \citep{nolle2012}. This topological coincidence suggests that Möbius-like operations (chromatic inversion with lightness continuity) may align with the intrinsic geometry of color space. The mathematical grounding for this connection is explored in detail in \S\ref{sec:mobius-mathematical}.

% ------------------------------------------------------------------------------
% 2.3 OKLab: Practical Implementation of Perceptual Uniformity
% ------------------------------------------------------------------------------
\subsection{OKLab Color Space}
\label{sec:oklab-implementation}

\subsubsection{Design Rationale}

Given the impossibility of perfect perceptual uniformity (\S\ref{sec:topology-impossibility}), the selection of a working color space requires balancing competing objectives:

\begin{enumerate}
    \item \textbf{Perceptual accuracy:} Euclidean distances should approximate perceptual differences as closely as possible
    \item \textbf{Computational efficiency:} Transformations should be fast enough for real-time generation
    \item \textbf{Numerical stability:} The space should behave predictably at gamut boundaries and extreme values
    \item \textbf{Industry validation:} The space should have independent verification and adoption
\end{enumerate}

OKLab, introduced by Björn Ottosson in 2020 \citep{ottosson2020}, achieves an excellent balance across these criteria. It has gained rapid adoption in web standards (CSS Color Level 4) \citep{csscolor4}, creative software, and independent validation for gradient quality \citep{levien2021oklab}.

Table~\ref{tab:color-space-comparison} synthesizes the trade-offs:

\begin{table}[htbp]
\centering
\caption{Color space comparison for temporal color journey construction}
\label{tab:color-space-comparison}
\begin{tabular}{lcccc}
\toprule
\textbf{Color Space} & \textbf{Uniformity} & \textbf{Cost} & \textbf{Stability} & \textbf{Selected} \\
\midrule
sRGB & Poor & Low & Good & No \\
HSV/HSL & Poor & Low & Poor & No \\
CIELAB (1976) & Moderate & Moderate & Poor at extremes & No \\
CAM16-UCS & Excellent & High & Good & No \\
\textbf{OKLab} & \textbf{Excellent} & \textbf{Low} & \textbf{Excellent} & \textbf{Yes} \\
\bottomrule
\end{tabular}
\end{table}

\textit{Uniformity} refers to how closely Euclidean distances match perceived differences (assessed via gradient smoothness and discrimination threshold alignment); \textit{Cost} measures computational complexity; \textit{Stability} indicates behavior at gamut boundaries and extreme lightness values. OKLab achieves CAM16-level uniformity with CIELAB-level computational cost.

\subsubsection{Coordinate Representation}

A color in OKLab is represented by three coordinates $(L, a, b)$:

\begin{description}
    \item[$L$ --- Perceived lightness] ranges from 0 (perceptual black) to 1 (brightest white). Equal differences in $L$ correspond to approximately equal perceived brightness changes.
    
    \item[$a$ --- Green–red opponent axis] with positive values indicating reddish tints and negative values indicating greenish tints. Zero is achromatic along this axis.
    
    \item[$b$ --- Blue–yellow opponent axis] with positive values indicating yellowish tints and negative values indicating bluish tints. Zero is achromatic along this axis.
\end{description}

The axes are designed to be \emph{approximately perceptually orthogonal}: altering $L$ alone should not significantly change perceived hue or chroma; altering $a$ or $b$ primarily affects hue and saturation without changing perceived lightness \citep{ottosson2020}.

The perceptual distance between two colors is computed as Euclidean distance:

\begin{equation}
    \Delta E = \sqrt{(L_2 - L_1)^2 + (a_2 - a_1)^2 + (b_2 - b_1)^2}
    \label{eq:oklab-distance}
\end{equation}

\noindent where $\Delta E$ denotes color difference in OKLab perceptual space.\footnote{All $\Delta E$ measurements in this specification refer to color difference calculated in OKLab space unless explicitly noted otherwise.}

\subsubsection{Transformation Pipeline}

The transformation from sRGB to OKLab follows a fixed matrix pipeline:

\begin{enumerate}
    \item \textbf{Linearize sRGB:} Remove gamma encoding to obtain linear RGB values
    \item \textbf{RGB to XYZ:} Apply the standard sRGB-to-XYZ matrix (D65 white point)
    \item \textbf{XYZ to LMS:} Apply matrix $M_1$ to approximate cone responses
    \item \textbf{Nonlinear compression:} Apply cube root: $(l', m', s') = (l^{1/3}, m^{1/3}, s^{1/3})$
    \item \textbf{LMS to Lab:} Apply matrix $M_2$ to obtain final $(L, a, b)$
\end{enumerate}

The transformation matrices are:

\begin{equation}
    M_1 = \begin{pmatrix}
    0.8189330101 & 0.3618667424 & -0.1288597137 \\
    0.0329845436 & 0.9293118715 & 0.0361456387 \\
    0.0482003018 & 0.2643662691 & 0.6338517070
    \end{pmatrix}
    \label{eq:oklab-m1}
\end{equation}

\begin{equation}
    M_2 = \begin{pmatrix}
    0.2104542553 & 0.7936177850 & -0.0040720468 \\
    1.9779984951 & -2.4285922050 & 0.4505937099 \\
    0.0259040371 & 0.7827717662 & -0.8086757660
    \end{pmatrix}
    \label{eq:oklab-m2}
\end{equation}

Each step models aspects of human vision: $M_1$ approximates L, M, S cone responses; the cube root models nonlinear perceptual compression (analogous to the logarithmic relationship described by Fechner \citeyearpar{fechner1860}); $M_2$ extracts lightness as a weighted average and opponent-color signals as cone differences.

\paragraph{Computational Efficiency}

The OKLab conversion is computationally trivial:
\begin{itemize}
    \item \textbf{Per conversion:} approximately 50–100 CPU cycles
    \item \textbf{Operations:} 18 multiplications, 12 additions, 3 cube roots
    \item \textbf{Memory:} no allocations; all register operations
    \item \textbf{Parallelization:} trivially vectorizable with SIMD
\end{itemize}

This efficiency enables real-time palette generation. The entire forward/inverse pipeline uses fixed, known constant matrices—no iterative solving, no conditionals, no special cases.

\subsubsection{Cartesian vs.\ Cylindrical Coordinates (OKLCh)}

For certain operations—particularly hue manipulation—a cylindrical form is more convenient. OKLCh represents the same colors using:

\begin{description}
    \item[$L$ --- Lightness] identical to OKLab
    \item[$C$ --- Chroma] defined as $C = \sqrt{a^2 + b^2}$, representing distance from the neutral axis
    \item[$h$ --- Hue angle] defined as $h = \operatorname{atan2}(b, a)$, the angle around the $a$–$b$ plane
\end{description}

The conversion is straightforward:
\begin{align}
    C &= \sqrt{a^2 + b^2} \label{eq:chroma} \\
    h &= \operatorname{atan2}(b, a) \label{eq:hue-angle}
\end{align}

The engine uses both representations internally:
\begin{itemize}
    \item \textbf{Cartesian (OKLab)} for distance calculations, linear interpolation, and Bézier curve construction
    \item \textbf{Cylindrical (OKLCh)} for hue-related operations such as warmth bias (\S\ref{sec:warmth-bias}), complementary color calculation (\S\ref{sec:mobius-loop}), and palette mode presets (\S\ref{sec:palette-modes})
\end{itemize}

This separation allows clean reasoning about lightness dynamics (along $L$) versus chromatic dynamics (in the $a$–$b$ plane or around $h$).

\subsubsection{Implications for Journey Construction}

The choice of OKLab has cascading implications for color journey design:

\begin{enumerate}
    \item \textbf{Linear interpolation produces smooth gradients:} A straight line between two colors in OKLab coordinates produces a perceptually smooth transition—unlike RGB, where intermediate steps may appear desaturated or shift unexpectedly in hue.
    
    \item \textbf{Distance constraints are meaningful:} The minimum step $\Delta_{\min} \approx 2.0$ and maximum step $\Delta_{\max} \approx 5.0$ correspond to actual perceptual differences (approximately 2× and 5× JND respectively), making these parameters interpretable.
    
    \item \textbf{Dynamics are approximately separable:} The orthogonality of $L$, $a$, $b$ means lightness can be adjusted independently of chromaticity—critical for implementing perceptual velocity weights (\S\ref{sec:velocity-weights}).
    
    \item \textbf{Hue interpolation is stable:} Unlike HSV, where interpolating across the 0°/360° boundary causes discontinuities, OKLab's Cartesian representation avoids such artifacts.
\end{enumerate}

Because OKLab is unbounded—it can represent colors outside any physical display gamut—the engine must include explicit gamut mapping (\S\ref{sec:gamut-management}).

% ------------------------------------------------------------------------------
% 2.4 Temporal vs. Spatial Color Perception
% ------------------------------------------------------------------------------
\subsection{Temporal vs.\ Spatial Color Perception}
\label{sec:temporal-spatial}

\subsubsection{Eye Movement and Chromatic Sensitivity}

Human color perception differs fundamentally between static spatial patterns and dynamic temporal transitions. Braun et al.\ \citeyearpar{braun2017} measured chromatic sensitivity during different types of eye movements:

\begin{description}
    \item[Smooth pursuit:] When tracking a moving object with smooth eye movements, chromatic sensitivity \emph{increases} by approximately 12\% \citep{braun2017}. The visual system enhances color discrimination during continuous tracking.
    
    \item[Saccadic suppression:] During rapid eye movements (saccades), chromatic sensitivity \emph{decreases} by approximately 58\% \citep{braun2017}. The visual system suppresses color information during ballistic gaze shifts to avoid motion blur.
\end{description}

This asymmetry has direct implications for color journey design: transitions should be optimized for \emph{smooth viewing} rather than rapid scanning. Gradual, continuous color changes exploit the enhanced chromatic sensitivity during pursuit, while abrupt changes risk perceptual suppression.

\subsubsection{Temporal Asymmetry in Color Channels}

Beyond eye movement effects, the visual system exhibits channel-specific temporal sensitivities. Sekulovski et al.\ \citeyearpar{sekulovski2007} measured smoothness thresholds—the maximum color difference between successive frames that observers still perceive as smooth—for lightness, chroma, and hue changes:

\begin{quote}
``The visibility threshold of smoothness, i.e., the maximum color difference between two successive colors that is allowed in order to perceive a temporal color transition as smooth, is about ten times smaller for lightness changes than for chroma or hue changes in CIELAB'' \citep[p.~530, citing][]{kong2021}; \citep{sekulovski2007}.
\end{quote}

In other words, observers are approximately \textbf{10 times more sensitive} to temporal changes in lightness than to changes in chromatic attributes \citep{sekulovski2007}. This 10:1 asymmetry is a robust empirical finding, replicated across base colors and viewing conditions.

Importantly, Sekulovski et al.\ \citeyearpar{sekulovski2007} found that lightness thresholds are \emph{independent} of base chromaticity:

\begin{quote}
``No main effect of the base color point on the thresholds for lightness changes was found for the chromatic color points\ldots\ This shows independence of the thresholds to changes in chromaticity and dependence on lightness'' \citep[Results section]{sekulovski2007}.
\end{quote}

This independence simplifies calibration: a single temporal sensitivity ratio can be applied globally.

\subsubsection{Design Implications}

These findings directly inform the engine's velocity weighting scheme (\S\ref{sec:velocity-weights}):

\begin{enumerate}
    \item \textbf{Lightness velocity weight:} $w_L = 10.0$ reflects the heightened temporal sensitivity \citep{sekulovski2007}
    \item \textbf{Chroma velocity weight:} $w_C = 1.0$ serves as baseline (unity weight)
    \item \textbf{Hue velocity weight:} $w_H \approx 1.5$ to $2.0$ is a design parameter requiring empirical validation
\end{enumerate}

The 10:1 ratio for $w_L : w_C$ is empirically validated \citep{sekulovski2007}. However, hue velocity weighting remains a \emph{design heuristic}; Sekulovski's data show hue and chroma thresholds are comparable but do not provide a precise ratio. Future work should employ psychometric methods (e.g., two-alternative forced-choice paradigms \citep{wang2022}) to measure hue temporal sensitivity rigorously.

\subsubsection{Kong's Gap: No Temporal Color Spaces}

Kong \citeyearpar{kong2021} provides a critical observation that frames the motivation for this work:

\begin{quote}
``CIELAB\ldots\ is not a useful space to predict the perception of dynamic colored light. Today, no color spaces are available that accurately predict the visibility of color differences over time'' \citep[p.~530]{kong2021}.
\end{quote}

Existing color spaces—CIELAB, CAM16, even OKLab—are calibrated for \emph{spatial} discrimination tasks: side-by-side color patches viewed simultaneously. They do not account for temporal dynamics: how quickly colors can change, how eye movements modulate perception, or how chromatic channels differ in temporal sensitivity.

The engine does not claim to solve this problem fully, but it \emph{acknowledges} the gap explicitly. By incorporating empirically-measured temporal asymmetries (Sekulovski's 10:1 ratio) and eye-movement effects (Braun's smooth pursuit enhancement), the design takes concrete steps toward temporally-aware color space usage.

% ------------------------------------------------------------------------------
% 2.5 Discrimination Thresholds and Just-Noticeable Differences
% ------------------------------------------------------------------------------
\subsection{Discrimination Thresholds and Just-Noticeable Differences}
\label{sec:jnd-thresholds}

\subsubsection{Empirical Measurement}

Hong et al.\ \citeyearpar{hong2024} provide comprehensive measurements of color discrimination thresholds across the visible gamut. Their study involved approximately 6,000 trials per participant, using a two-alternative forced-choice (2AFC) design to map elliptical contours of equal discriminability.

The just-noticeable difference (JND)—the smallest color difference a typical observer can reliably detect—varies across color space due to its Riemannian structure (\S\ref{sec:riemannian-color}). However, modern uniform spaces like OKLab are designed such that $\Delta E \approx 1.0$ to $2.0$ approximates the JND for most observers under standard viewing conditions \citep{hong2024,ottosson2020}.

\subsubsection{Practical JND Threshold}

The engine adopts $\Delta E \approx 2.0$ as a conservative practical JND threshold:

\begin{itemize}
    \item \textbf{Minimum step:} $\Delta_{\min} = 2.0$ ensures successive colors are perceptually distinguishable
    \item \textbf{Maximum step:} $\Delta_{\max} = 5.0$ is a design heuristic (approximately 2–3× JND) chosen to maintain perceptual continuity
\end{itemize}

The $\Delta_{\min}$ threshold is grounded in empirical discrimination data \citep{hong2024}. The $\Delta_{\max}$ threshold, however, lacks direct empirical validation for temporal transitions. It represents an engineering choice informed by the principle that ``small multiples of JND'' should preserve perceptual coherence. Future research should measure closure acceptability thresholds using psychometric methodology \citep{wang2022}.

\subsubsection{Temporal vs.\ Spatial JND}

Importantly, the JND values from Hong et al.\ \citeyearpar{hong2024} and Ottosson \citeyearpar{ottosson2020} are measured for \emph{simultaneous spatial comparisons} (side-by-side patches). Temporal discrimination—how different two colors must be to notice a change when viewing them sequentially—may differ. Sekulovski et al.\ \citeyearpar{sekulovski2007} measured temporal smoothness thresholds but did not provide a single universal JND equivalent.

This uncertainty is acknowledged explicitly in the specification:

\begin{quote}
While the 10:1 lightness-chroma asymmetry is empirically validated \citep{sekulovski2007}, threshold parameters ($\Delta_{\min}$, $\Delta_{\max}$) for temporal transitions remain design heuristics requiring future empirical validation.
\end{quote}

% ------------------------------------------------------------------------------
% 2.6 Limitations and Design Heuristics
% ------------------------------------------------------------------------------
\subsection{Limitations and Design Heuristics}
\label{sec:limitations-heuristics}

\subsubsection{Validated Parameters}

The following design parameters are grounded in empirical research:

\begin{description}
    \item[Riemannian manifold structure:] Confirmed by comprehensive discrimination measurements \citep{hong2024}
    \item[720° hue topology:] Mathematically derived from metric tensor measurements \citep{nolle2012}
    \item[Temporal lightness-chroma asymmetry (10:1):] Empirically measured via smoothness threshold experiments \citep{sekulovski2007}
    \item[JND approximation ($\Delta E \approx 1$–$2$):] Validated for spatial discrimination in OKLab \citep{hong2024,ottosson2020}
    \item[Smooth pursuit chromatic enhancement (+12\%):] Measured via 2AFC sensitivity tasks during eye tracking \citep{braun2017}
\end{description}

\subsubsection{Design Heuristics Requiring Validation}

The following parameters are \emph{design choices} informed by the literature but lacking direct empirical validation:

\begin{description}
    \item[Maximum step threshold ($\Delta_{\max} = 5.0$):] Chosen as 2–3× spatial JND; no published validation for temporal closure perception
    \item[Hue velocity weight ($w_H \approx 1.5$ to $2.0$):] Inferred from Sekulovski's data showing hue and chroma have comparable sensitivity; precise ratio unmeasured
    \item[Smoothstep easing function:] Selected for C² continuity; no perceptual validation comparing easing curves for color transitions
    \item[Closure threshold ($\Delta_{\text{close}} \leq \Delta_{\max}$):] Derived from coherence reasoning; not empirically tested for return-to-origin acceptability
\end{description}

\subsubsection{Recommended Validation Methodology}

Future research should employ psychometric methods to validate these heuristics. Wang et al.\ \citeyearpar{wang2022} demonstrate the use of two-alternative forced-choice (2AFC) paradigms with Point of Subjective Equality (PSE) analysis for measuring perceptual thresholds in dynamic contexts. Specific studies recommended:

\begin{enumerate}
    \item \textbf{Closure threshold detection:} Present observers with color loops of varying $\Delta_{\text{close}}$ and measure detection rates for ``returns to origin'' vs.\ ``does not return''
    \item \textbf{Hue velocity weight calibration:} Measure temporal smoothness thresholds for pure hue changes and fit $w_H$ to match lightness-chroma asymmetry data
    \item \textbf{Maximum step acceptability:} Measure subjective quality ratings for journeys with varying $\Delta_{\max}$ to determine perceptual coherence boundaries
\end{enumerate}

\subsubsection{Acknowledgment of Scope}

This specification grounds perceptual claims in published research wherever possible \citep{hong2024,sekulovski2007,braun2017,nolle2012}. Where empirical evidence is unavailable, parameters are framed explicitly as design heuristics subject to future validation. This transparency is essential for academic integrity and distinguishes evidence-based choices from engineering conveniences.

% ==============================================================================
% End of Section 2
% ==============================================================================
